\section{Conclusiones}

El objetivo principal de \textbf{8-Track} se cumplió de manera satisfactoria. Se lograron implementar los cuatro
módulos principales: \textbf{minero}, \textbf{base de datos}, \textbf{compilador}, \textbf{interfaz gráfica}, y
funcionaron de manera conjunta. Además, se implementó la reproducción de audio como una extensión del sistema
y se logró un flujo de aplicación usable.

Aún asi, me encontré con varios problemas durante el desarrollo del proyecto. En un inicio me fue bastante difícil
abstraer correctamente el planteamiento (a nivel código) de la interfaz gráfica, por eso terminé con bastante código
repetido entre los componentes \texttt{QML}. El desarrollo del \texttt{parser} de \textbf{MuQL} fue bastante complicado,
y me tomó días de lectura asimilarlo.

Como trabajo futuro, se plantea ampliar la compatibilidad con otros formatos de audio, incorporar un modo de reproducción
más completo e inmersivo e integrar APIs externas que permitan enriquecer y complementar los metadatos de los archivos.
